%=============================================================================!
% 0.1  THE QUESTION THIS BOOK PREPARES FOR                                    !
%=============================================================================!
\section{The question this book prepares for}
\label{sec:or-question}

Molecular dynamics follows atoms by integrating Newton's equations in small
steps of time. The forces may come from a prescribed function of the
atomic positions, an electronic-structure calculation or a function fitted
by machine learning. In each case, the step must resolve the fastest
vibration retained in the model; Chapter~11 explains how fixing selected
bond lengths removes some of these vibrations. In a graphite anode with a
carbonate electrolyte, the stretches of carbon--hydrogen and carbon--carbon
bonds hold it near a femtosecond, while lithium moving between sites,
solvent reorganising around an ion and an interphase growing at the surface
are many orders of magnitude slower.

Machine-learned potentials can reduce the cost of each force evaluation,
but a cheaper evaluation does not by itself permit a longer stable step.
Learned propagators predict the motion over a longer interval. TrajCast maps
the positions and velocities at time $t$ to those at $t + \Dt$, with $\Dt$
ten to thirty times the reference step, and repeats the map to build a
trajectory\cite{Thiemann2026}; ATOM predicts several future configurations
in one evaluation and is pretrained across many molecules\cite{Thompson2026}.
TrajCast was tested on a molecule, crystalline quartz and liquid water;
ATOM was tested on trajectories of isolated molecules. Whether such
models can describe an electrochemical system, with its periodic electrode,
mobile ions, long-range electrostatics and rare activated events, is the
question this book prepares its reader to investigate.

Answering it draws on four bodies of knowledge, each resting on the one
before: classical mechanics and its numerical integration, which define a
trajectory and limit its step; statistical mechanics, which ties a
trajectory to the ensemble it samples and explains what a thermostat does
to it; the practice of molecular dynamics, from converged averages to the
observables both papers use as evidence; and machine learning, as far as the
equivariant networks and transformers on which the two models are built.
The present edition develops the first three and begins the fourth with
regression and optimisation; the course map identifies the planned
continuation.
